\documentclass[a4paper,11pt,twocolumn]{article}

% --- Encoding & Swedish language support ---
\usepackage[utf8]{inputenc}      % Allows å, ä, ö etc. directly in the source file
\usepackage[T1]{fontenc}         % Proper font encoding for accented/special characters
\usepackage{lmodern}             % Scalable fonts that match T1 encoding (crisper output)
\usepackage[swedish]{babel}      % Swedish hyphenation, date formats, and typographic conventions

% --- Layout & formatting ---
\usepackage[margin=2.5cm]{geometry}
\usepackage{enumitem}            % Nicer control over lists
\usepackage{titlesec}            % Custom section headings
\usepackage{parskip}             % Space between paragraphs instead of indentation

\titleformat{\section}{\Large\bfseries}{\thesection}{1em}{}
\titleformat{\subsection}{\large\bfseries}{\thesubsection}{1em}{}

\title{Recept}
\author{Gunnar Lindholm}
\date{\today}

\begin{document}

\maketitle
\thispagestyle{empty}

\clearpage
\begingroup
\small
\tableofcontents
\thispagestyle{empty}
\endgroup
\clearpage

\section{Allmänna råd - viktigt!}
\begin{itemize}
    \item Skaffa en bra assistent som kan knåda degen. Det blir bättre än att göra för hand.
    \item När det i bröd recept står "Gör som vanligt" så betyder det:
    \begin{enumerate}
        \item Smula jäst i bunken
        \item Värm degspadet till max 37 grader. Detta kan göras genom att bara värma vätska, eller om det ingår margarin, smälta detta och hälla på vätskan. Det kräver lite erfarenhet att inte få det för varmt. Jästen \emph{dör} om den utsätts för vätska som är för varm.
        \item Se till att jästen löser sig i vätskan. Det kan vara enklast att hälla i lite av vätskan först.
        \item Häll i all vätska, kryddor, salt, m.m.
        \item Häll i mjöl, lite i taget medan assistenten jobbar eller du rör för hand. Låt gärna assistenten jobba länge med degen.
        \item När allt mjöl är i, arbeta degen ordentligt.
        \item Låt degen jäsa under bakduk (ej i drag så stäng fönster)
    \end{enumerate}
    \item Ta lite extra mjöl på bakbordet när du bakar ut efter första jäsningen.
    \item När bröd är färdiggräddat så ska det låta lite ihåligt när man knackar på det på skorpan undertill.
    \item När bröd är färdig gräddat så är det bäst att låta det svalna på galler under bakduk.
    \item Om du ska ha ugn på X grader så sätter du en varmluftsugn på ungefär 10\% mindre grader.
    \item Kyckling måste tillagas ordentligt. Den ska vara vit inuti och saften genomskinlig. Observera att den kan bli torr om den tillagas för länge.
\end{itemize}

\section{Bröd - Frallor}




\subsection{Brytbröd}
\subsubsection*{Ingredienser}
\begin{itemize}[noitemsep]
    \item 50 g margarin
    \item 5 dl mjölk
    \item 50 g jäst
    \item 1 tsk salt (ev. mer)
    \item 1 msk socker
    \item 750 g vetemjöl
\end{itemize}

\subsubsection*{Gör så här}
\begin{enumerate}[label=\arabic*.]
    \item Gör som vanligt.
    \item Jäs 30 min.
    \item Baka ut 32 runda bullar (de blir små) och lägg dem tätt i hexagon mönster.
    \item Jäs 30 min.
    \item Pensla uppvispat ägg och strö på sesamfrö eller mörka vallmofrön.
    \item Grädda mitt i ugn i 225 grader i 20 min. Lägg på folie efter ungefär halva tiden ifall de ser ut att bli brända.
    Undersök i framtiden om detta recept kan bli semelbullar vid ökande av socker och tillsättande av kardemumma.
\end{enumerate}


\subsection{Frukostbröd}
Detta jäser över natten och gräddas på morgonen.
\subsubsection*{Ingredienser}
\begin{itemize}[noitemsep]
    \item 10g jäst
    \item 15g salt
    \item 300g vatten (kallt)
    \item 350g vetemjöl
    \item 50g grovt rågmjöl
\end{itemize}
\subsubsection*{Gör så här}
\begin{enumerate}[label=\arabic*.]
    \item Gör som vanligt. Kör 10 minuter i degblandare.
    \item Vila i 30 min.
    \item Dela i 8 delar och forma runda.
    \item Lägg en fuktig handduk på t.ex. en bricka och lägg degbitarna ovanpå
    \item Mjöla lätt och lägg en fuktig tunn handduk ovanpå.
    \item Låt jäsa 10-14 timmar över natten.
    \item Lägg på en plåt med papper och snitta
    \item Spraya ugnsbotten med vatten. Grädda i 10-15 min i 250 grader, sedan 5 min i 220 grader.
\end{enumerate}


\subsection{Gunnars frallor (småfranska)}
\subsubsection*{Ingredienser}
\begin{itemize}[noitemsep]
    \item 100g jäst
    \item 13 dl vatten
    \item 1,3 dl matolja
    \item 800g rågsikt
    \item 1440g vetemjöl
    \item 1,5 msk salt
\end{itemize}
\subsubsection*{Gör så här}
\begin{enumerate}[label=\arabic*.]
    \item Gör som vanligt.
    \item Jäs 40 min.
    \item Baka ut 64 frallor.
    \item Jäs 30 min.
    \item Grädda ca 12 min i 200 grader.
\end{enumerate}


\subsection{Hamburgerbröd / Korvbröd}
32 Hamburgerbröd eller 48 korvbröd. Godast samma dag. Går bra att frysa.
\subsubsection*{Ingredienser}
\begin{itemize}[noitemsep]
    \item 100g jäst
    \item 100g margarin
    \item 10 dl mjölk
    \item 2 tsk salt
    \item 1 msk socker
    \item 1,5 kg vetemjöl
    \item 1-2 ägg (till Hamburgerbröd)
    \item sesamfrön (till Hamburgerbröd)
\end{itemize}
\subsubsection*{Gör så här}
\begin{enumerate}[label=\arabic*.]
    \item Gör som vanligt.
    \item Baka länge och rejält i bakmaskinen.
    \item Jäs 45 minuter
    \item Häll ut degen och dela genast till rätt antal bitar (ej knåda)
    \item Platta till hamburgerbröden
    \item Jäs i 20 min.
    \item Pensla hamburgerbröden med ägg och strö över sesamfrön.
    \item Grädda i 225 grader i 8-10 minuter.
\end{enumerate}



\newpage
\section{Bröd - Limpor}


\subsection{Baguetter}
\subsubsection*{Ingredienser}
\begin{itemize}[noitemsep]
    \item 75g jäst
    \item 5 dl vatten
    \item 2 tsk salt
    \item 810g vetemjöl (ev. mer)
\end{itemize}
\subsubsection*{Gör så här}
\begin{enumerate}[label=\arabic*.]
    \item Gör degen som vanligt
    \item Jäs 1 h.
    \item Knåda på bakbord (ev. mer mjöl)
    \item Dela i 4 delar, kavla ut och rulla ihop (eller bara gör längder)
    \item Skåra dem.
    \item Jäs 15-20 min.
    \item Pensla kallt vatten.
    \item 275 grader ca 15 minuter.
    \item Svalna utan bakduk = knaprig skorpa, med bakduk = mjukare skorpa.
\end{enumerate}



\subsection{Gotlandsknutar}
\subsubsection*{Ingredienser}
\begin{itemize}[noitemsep]
    \item 75 g jäst
    \item 5 dl mjölk
    \item 1 dl olja
    \item 0,5 dl sirap
    \item 2 tsk salt
    \item 1 msk pomeransskal
    \item 1 tsk stött anis
    \item 550 g rågsikt
    \item 420 g vetemjöl
\end{itemize}
\subsubsection*{Gör så här}
\begin{enumerate}[label=\arabic*.]
    \item Gör som vanligt (men ej vetemjöl)
    \item Vila i 10 minuter
    \item Tillsätt vetemjöl lite i sänder och arbeta smidig.
    \item Låt ej jäsa, utan dela i 3 delar på bakbord.
    \item Rulla till 1 m långa.
    \item Forma till en "knut". T.ex. genom att låta höger ända vara still, forma en ögla (så att vänster änden sticker ut åt vänster och kommer ut underifrån) genom att med höger hand vriden moturs så att tummen pekar åt höger (obekvämt) ta tag i "korven" ovanifrån, lyft och vrid tillbaka handen så att en ögla bildas. Stoppa sedan ner den vänstra änden i öglan och sammanfoga den med den högra änden, så att skarven hamnar under en del av öglan.
    \item Låt jäsa 30 minuter.
    \item Pensla med kallt vatten och strö på vetemjöl.
    \item Grädda i 200 grader i 20 minuter.
\end{enumerate}



\subsection{Greta Kyrks Lingonbröd}
\subsubsection*{Ingredienser}
\begin{itemize}[noitemsep]
    \item 4 dl ljummet vatten
    \item 1 dl vetekross
    \item 1 dl rågkross
    \item 1 msk salt
    \item 1 msk kummin
    \item 2 dl lingonsylt
    \item 1 dl russin
    \item 1 dl linfrön
    \item 1 dl vetekli
\end{itemize}
Blanda allt ovan och låt stå över natten
\begin{itemize}[noitemsep]
    \item 2 pkt torrjäst
    \item 1,2 kg vetemjöl
    \item 5 dl kokande vatten
\end{itemize}
\subsubsection*{Gör så här}
\begin{enumerate}[label=\arabic*.]
    \item Blanda jästen väl med mjölet
    \item Häll kokande vatten över gårdagens blandning.
    \item Blanda sedan i mjöl och jäst blandningen.
    \item Degen behöver ej knådas.
    \item Forma genast två limpor som läggs i smorda 2 L formar som ska jäsa snabbt i 50 grader i 10 min.
    \item Jäs sedan utanför ugnen till dubbel storlek (ca 20 min).
    \item Grädda nederst i 250 grader i 30 min, samt 10 min i eftervärme. Använd folie om brända.
    \item Svalna.
\end{enumerate}



\subsection{Grovt formbröd}
\subsubsection*{Ingredienser}
\begin{itemize}[noitemsep]
    \item 50 g jäst
    \item 2 tsk salt
    \item 1 msk anis
    \item 1 msk fänkål
    \item 1 dl sirap
    \item 9 dl fint rågmjöl
    \item 5 dl vatten
    \item 4-5 dl vetemjöl
\end{itemize}
\subsubsection*{Gör så här}
\begin{enumerate}[label=\arabic*.]
    \item Gör som vanligt
    \item Häll på lite mjöl
    \item Jäs 30 min.
    \item Smörj 2 formar.
    \item Baka ut degen i formarna.
    \item Jäs x minuter
    \item Grädda i 200 grader, 35-40 minuter
\end{enumerate}



\subsection{Halltorpsvört}
\subsubsection*{Ingredienser}
\begin{itemize}[noitemsep]
    \item 4 dl svagdricka + 66 cl porter
    \item 100g jäst
    \item 1/2 dl olja
    \item 1 msk salt
    \item 100g russin
    \item 1,5 tsk ingefärra
    \item 1,5 tsk kryddnejlikor
    \item 1 msk malda pomeransskal
    \item 1 dl sirap
    \item 2,2 L rågsikt
    \item 7 dl vetemjöl
\end{itemize}
\subsubsection*{Gör så här}
\begin{enumerate}[label=\arabic*.]
    \item Gör som vanligt
    \item Jäs 30 min.
    \item Dela till 3 limpor. Pensla med lite olja mellan limporna. Ev. skåra.
    \item Jäs x minuter
    \item Grädda nederst i 40 min.
\end{enumerate}




\subsection{Hålkakor}
\subsubsection*{Ingredienser}
\begin{itemize}[noitemsep]
    \item 50 g margarin
    \item 5 dl vatten
    \item 50 g jäst
    \item 2 tsk salt
    \item 0,5 dl sirap
    \item 2 tsk anis
    \item 2 tsk fänkål
    \item 550-715 g rågsikt
    \item 3 dl Finax fullkornsmjöl (eller annat)
\end{itemize}
\subsubsection*{Gör så här}
\begin{enumerate}[label=\arabic*.]
    \item Gör som vanligt
    \item Jäs 40 min.
    \item Kavla 4 kakor
    \item Gör hål i mitten med fingrarna.
    \item Lägg 2 och 2 på plåtar
    \item Nagga.
    \item Jäs 30 min.
    \item Grädda 225 grader i 15-20 min.
    \item Ev. pensla med vatten efter gräddning.
\end{enumerate}



\subsection{Korgbröd}
\subsubsection*{Ingredienser}
\begin{itemize}[noitemsep]
    \item 50 g jäst
    \item 50 g margarin
    \item 4,5 dl vatten
    \item 1 msk salt
    \item 1 msk pomeransskal
    \item 3 msk apelsinmarmelad
    \item 605 g rågsikt
    \item 150 g vetemjöl
\end{itemize}
\subsubsection*{Gör så här}
\begin{enumerate}[label=\arabic*.]
    \item Gör som vanligt. Ev. mer mjöl vid behov.
    \item Jäs 40 minuter
    \item Knåda till en stor bulle
    \item Mjöla en korg (mycket trillar ut)
    \item Jäs i korgen, 30 min.
    \item Vänd snabbt korgen så brödet trillar ut på plåten.
    \item Grädda längst ner i 200 grader i 30 min.
\end{enumerate}


\subsection{Kroppkakebröd}
\subsubsection*{Ingredienser: Dag 1}
\begin{itemize}[noitemsep]
    \item 1,5 L kroppkakspad
    \item 1,1 kg grovt rågmjöl
    \item 2,5 msk kummin
\end{itemize}
\subsubsection*{Gör så här: Dag 1}
\begin{enumerate}[label=\arabic*.]
    \item Blanda Ingredienserna väl.
    \item Värm spadet till 70 grader och häll över det torra.
    \item Täck till nästa dag.
\end{enumerate}
\subsubsection*{Ingredienser: Dag 2}
\begin{itemize}[noitemsep]
    \item 50g jäst
    \item 1 dl mörk sirap
    \item 1,5 dl surdeg
    \item 4 tsk salt
    \item 1,2 kg rågsikt
\end{itemize}
\subsubsection*{Gör så här: Dag 2}
\begin{enumerate}[label=\arabic*.]
    \item Lös upp jästen med lite vatten.
    \item Blanda i detta i degblandningen från igår.
    \item Tillsätt sedan sirap, salt och rågmjöl (lite i sänder, det blir tungt). Ev mer mjöl.
    \item Ta upp, knåda väl
    \item Forma 3 limpor i smord långpanna. Pensla olja mellan dem
    \item Jäs 40-60 minuter
    \item Nagga och grädda långt ner i 200 grader ca 1h.
    \item Svalna i bakduk.
    \item Låt vila till nästa dag.
\end{enumerate}



\subsection{Moster Hulda}
\subsubsection*{Ingredienser}
\begin{itemize}[noitemsep]
    \item 13 dl vatten
    \item 2 msk salt
    \item 2 dl sirap
    \item 50 g jäst
    \item 3 msk kummin
    \item 2 kg rågsikt
    \item 2-3 dl solrosfrön \emph{eller} 3 dl vetekross som fått svälla en stund i lite vatten.
\end{itemize}
\subsubsection*{Gör så här}
\begin{enumerate}[label=\arabic*.]
    \item Gör som vanligt
    \item Gör 3 limpor
    \item Jäs 30 min på plåt
    \item Skåra brödet
    \item Grädda 30 min i 225 grader
    \item Stäng av ugnen men låt brödet stå kvar under natten.
\end{enumerate}



\subsection{Sötebröd från Medelpad}
\subsubsection*{Ingredienser}
\begin{itemize}[noitemsep]
    \item 2,5 L vatten
    \item 4,2 kg rågmjöl
    \item 5 dl sirap
    \item 1 msk fänkål
    \item 1 msk anis
    \item 100g jäst
\end{itemize}
\subsubsection*{Gör så här - dag 1}
\begin{enumerate}[label=\arabic*.]
    \item Koka vattnet
    \item Blanda med mjöl i en bunke, rör om noga
    \item Låt stå varmt i ett dygn.
\end{enumerate}
\subsubsection*{Gör så här - dag 2}
\begin{enumerate}[label=\arabic*.]
    \item Rör ut jästen i lite ljummet vatten.
    \item Häll det och kryddor och sirap över gårdagens blandning.
    \item Låt jäsa upp.
    \item Knåda (ev lite vetemjöl)
    \item Låt jäsa igen.
    \item Dela i två delar.
    \item Jäs på plåtar.
    \item Grädda i 150-175 grader i 2h.
\end{enumerate}


\subsection{Ungerskt lantbröd}
\subsubsection*{Ingredienser}
\begin{itemize}[noitemsep]
    \item 50 g jäst
    \item 25 g margarin
    \item 7,5 dl vatten
    \item 1 msk salt
    \item 330g rågsikt
    \item 850g vetemjöl (ofta lite mer)
    \item Ev. 1-1,5 msk kummin.
\end{itemize}
\subsubsection*{Gör så här}
\begin{enumerate}[label=\arabic*.]
    \item Gör som vanligt. Arbeta degen länge.
    \item Jäs 1,5 h.
    \item Forma två stora limpor.
    \item Jäs 45 min.
    \item Skåra och grädda nederst i 225 grader i 40-50 min.
\end{enumerate}



\subsection{Zopf}
\subsubsection*{Ingredienser}
\begin{itemize}[noitemsep]
    \item 50 g jäst
    \item 5 dl mjölk
    \item 2 tsk salt
    \item 50 g margarin
    \item 1 kg vetemjöl
\end{itemize}
\subsubsection*{Gör så här}
\begin{enumerate}[label=\arabic*.]
    \item Gör som vanligt.
    \item Jäs 30 min.
    \item Dela i 7 delar (eller gör 3 delar i dessa proportioner)
    \begin{itemize}
        \item 4 delar som flätas och läggs i botten
        \item 2 delar som flätas och läggs ovanpå
        \item 1 del som flätas och läggs överst
    \end{itemize}
    Det är viktigast att understa delen blir stabil så det inte "välter".
    \item Peta in ändarna under den understa flätan.
    \item Jäs i 20 min.
    \item Pensla med ägg, strö på valmofrön.
    \item Grädda nederst i 225 grader ca 20 min.
\end{enumerate}


\newpage
\section{Knäckebröd}

\subsection{Goda knäckebröd}
\subsubsection*{Ingredienser}
\begin{itemize}[noitemsep]
    \item 2 dl majsmjöl
    \item 1 dl sesamfrön
    \item 1 dl solroskärnor
    \item 1 dl pumpakärnor
    \item 1/2 dl linfrön
    \item 1/2 dl rapsolja
    \item 2,5 dl kokande vatten
\end{itemize}
\subsubsection*{Gör så här}
\begin{enumerate}[label=\arabic*.]
    \item Blanda allt torrt
    \item Slå på olja och vatten
    \item Blanda och kavla ut på en plåt mellan två bakplåtspapper. Eller 1,5 plåt för att få extra tunnt.
    \item Grädda i 150 grader i ca 1h. Det ska vara hårt när man känner på det.
\end{enumerate}




\newpage
\section{Bakverk}




\subsection{Äppelkaka}
\subsubsection*{Ingredienser}
\begin{itemize}[noitemsep]
    \item 2 ägg
    \item 2 dl socker
    \item 2 3/4 dl vetemjöl
    \item 2 tsk bakpulver
    \item 2 tsk vaniljsocker
    \item 50 g smält smör
    \item 3/4 dl mjölk
    \item 2-3 äpplen i tunna klyftor
    \item kanel
\end{itemize}
\subsubsection*{Gör så här}
\begin{enumerate}[label=\arabic*.]
    \item Sätt ugnen på 200 grader
    \item Smörj och bröa en tårtform 22-25 cm diameter
    \item Vispa ägg och socker poröst
    \item Blanda mjöl, bakpulver och vaniljsocker.
    \item Tillsätt det smälta smöret, blandat med mjölken.
    \item Rör om till jämn smet
    \item Häll 1/3 av smeten i formen.
    \item Fördela äppelklyftorna jämt i formen och strö lite kanel på dem.
    \item Häll i resten av smeten (alla äppelbitar behöver inte bli täckta)
    \item Grädda i 175-200 grader omkring 35-60 minuter. Folie om bränd.
\end{enumerate}




\newpage
\section{Paj}




\newpage
\section{Såser}

\newpage
\section{Huvudrätter - soppa}
\subsection{Gunnars goda soppa}
Mycket uppskattad grönsak/köttsoppa. Det här blir storkok som kräver c:a 15 Liter grytkapacitet och räcker i veckor så du vill nog skala ner receptet
\subsubsection*{Ingredienser}
\begin{itemize}[noitemsep]
    \item 5 kg potatis
    \item 1,3 kg lök
    \item 2 kg köttfärs
    \item 2 kg morötter
    \item 8 buljongtärningar (kött och grönsak)
    \item 1-2 dl tomatpure (eller kethcup)
    \item 1-2 dl grädde (eller mer)
    \item salt
    \item kryddor: t.ex. vit- och svartpeppar, vitlök, oregano, salvia, franska örtkryddor
\end{itemize}
\subsubsection*{Gör så här}
\begin{enumerate}[label=\arabic*.]
    \item Skala potatisen och skiva och börja koka
    \item Hacka löken och släng i allt eftersom
    \item Bryn köttfärsen och släng i allt eftersom
    \item Ha i buljongtärningar
    \item Krydda och ha i tomatpure och grädde
    \item Låt det koka och smaka av
    \item Grov riv morötterna och häll i.
    \item Låt det puttra några minuter så morötterna blir mjuka.
\end{enumerate}


\newpage
\section{Huvudrätter - bröd}

\subsection{Pizza sås}
\subsubsection*{Ingredienser}
\begin{itemize}[noitemsep]
    \item Passerade tomater
    \item Salt
    \item Vitlök
    \item Oregano
\end{itemize}
\subsubsection*{Gör så här}
\begin{enumerate}[label=\arabic*.]
    \item Koka och smaka av.
\end{enumerate}


\subsection{Fattiga riddare}
\subsubsection*{Ingredienser}
\begin{itemize}[noitemsep]
    \item 2 ägg
    \item 1,5 dl mjölk
    \item 1 tsk socker
    \item margarin till stekning
    \item 8 skivor vittbröd
\end{itemize}
\subsubsection*{Gör så här}
\begin{enumerate}[label=\arabic*.]
    \item Vispa samman allt
    \item Doppa skivorna i äggsmeten och stek i pannan tills gyllenbruna
Servera med äppelmos eller lingon.
\end{enumerate}



\section{Huvudrätter - kyckling}

\subsection{Snabb enkel kyckling i tikka masala sås}
\subsubsection*{Ingredienser}
\begin{itemize}[noitemsep]
    \item 550 g strimlad kyckling
    \item 1 burk Tikka Masala sås
    \item 1 burk kokosmjölk (eventuellt)
    \item salt
    \item vitpepper
    \item Eventuellt lök eller andra grönsaker.
\end{itemize}
\subsubsection*{Gör så här}
\begin{enumerate}[label=\arabic*.]
    \item Fräs ev. lök
    \item Fräs strimlade kycklingen
    \item Salta och peppra lite grand
    \item Häll i Tikka Masala såsen
    \item Häll i kokosmjölk om så önskas (blir mildare men utan smak från kokosmjölken)
    \item Låt koka så kycklingen blir genomkokt. Det går fort.
    \item Häll ev. andra grönsaker.
\end{enumerate}


\subsection{Snabb enkel kyckling i vit sås}
\subsubsection*{Ingredienser}
\begin{itemize}[noitemsep]
    \item 550 g strimlad kyckling
    \item 1 burk kokosmjölk
    \item 1 dl Santa Maria Sweet Chilli sås
    \item salt
    \item vitpepper
    \item Eventuellt lök eller andra grönsaker.
\end{itemize}
\subsubsection*{Gör så här}
\begin{enumerate}[label=\arabic*.]
    \item Fräs ev. lök
    \item Fräs strimlade kycklingen
    \item Salta och peppra lite grand
    \item Häll i kokosmjölk och sweet chilli sås (filtrera bort fasta delar genom en sil, ev. ta mer sås)
    \item Låt koka så kycklingen blir genomkokt. Det går fort.
    \item Häll ev. andra grönsaker.
\end{enumerate}





\newpage
\section{Huvudrätter - kött}

\subsection{Köttbullar med Gräddsås}

Klassiska svenska köttbullar serverade med gräddsås, potatismos och lingonsylt.

\subsubsection*{Ingredienser}
\begin{itemize}[noitemsep]
    \item 500 g blandfärs (nöt och fläsk)
    \item 1 gul lök, finhackad
    \item 1 ägg
    \item 1 dl ströbröd
    \item 1 dl mjölk
    \item Salt och peppar efter smak
    \item 1 msk smör till stekning
    \item Till gräddsåsen: 3 dl grädde, 2 msk soja, 1 msk maizena
\end{itemize}

\subsubsection*{Gör så här}
\begin{enumerate}[label=\arabic*.]
    \item Blötlägg ströbrödet i mjölken i ca 10 minuter.
    \item Fräs den finhackade löken i smör tills den är mjuk, låt svalna.
    \item Blanda färs, ägg, ströbrödsblandning och lök. Krydda med salt och peppar.
    \item Rulla blandningen till jämnstora bullar med hjälp av kalla händer.
    \item Stek köttbullarna i smör på medelvärme tills de är genomstekta och fint bruna.
    \item Häll bort överflödigt fett, tillsätt grädde och soja i samma stekpanna.
    \item Red av såsen med utrört maizena och låt sjuda några minuter.
    \item Servera köttbullarna med såsen, kokt potatis eller mos och lingonsylt.
\end{enumerate}

\newpage
\section{Huvudrätter - potatis}

\subsection{Kroppkakor - öländska}
\subsubsection*{Ingredienser}
\begin{itemize}[noitemsep]
    \item $x$ kg potatis
    \item gul lök
    \item rimmat sidfläsk utan svål
    \item salt
    \item vitpeppar
    \item kryddpeppar
\end{itemize}
\subsubsection*{Gör så här}
\begin{enumerate}[label=\arabic*.]
    \item Koka och pressa potatis
    \item Riv rå potatis och pressa ut vatten genom duk (häll inte det pressade i vasken)
    \item Blanda potatisen till en smet (3/4 rå pressad)
    \item Tärna fläsket (lättare att skära smått om det är lite fruset)
    \item Hacka lök fint.
    \item Blanda fläsk, lök, salt och peppar (ev. stek lite för att smaka om kryddningen blev rätt)
    \item Gör bollar, platta ut och lägg på en klick med fyllning
    \item Slut bollen och se till att det inte är sprickor i den eller fyllning som sticker ut.
    \item Lägg ner alla samtidigt (dock en och en försiktigt med hålslev) i kokande vatten. Om det ej kokar så löses de upp. Ibland får man vänta med att lägga ner nästa.
    \item Häll i salt när alla är i 2-3 msk
    \item Koka 1 timme. Ta upp med hålslev. Gärna en i trä så man inte skadar dem med metallkant.
    \item Gör bröd av vattnet de kokat i.
\end{enumerate}


\newpage
\section{Efterrätter}





\subsection{x}
\subsubsection*{Ingredienser}
\begin{itemize}[noitemsep]
    \item
    \item
    \item
    \item
    \item
\end{itemize}
\subsubsection*{Gör så här}
\begin{enumerate}[label=\arabic*.]
    \item
    \item
    \item
    \item
    \item
    \item
    \item
    \item
    \item
\end{enumerate}

Tårta
Lufsa
Raggmunk
Grå kroppkakor
Vita kroppkakor

\end{document}
